% Options for packages loaded elsewhere
\PassOptionsToPackage{unicode}{hyperref}
\PassOptionsToPackage{hyphens}{url}
\PassOptionsToPackage{dvipsnames,svgnames,x11names}{xcolor}
%
\documentclass[
  letterpaper,
  DIV=11,
  numbers=noendperiod]{scrartcl}

\usepackage{amsmath,amssymb}
\usepackage{iftex}
\ifPDFTeX
  \usepackage[T1]{fontenc}
  \usepackage[utf8]{inputenc}
  \usepackage{textcomp} % provide euro and other symbols
\else % if luatex or xetex
  \usepackage{unicode-math}
  \defaultfontfeatures{Scale=MatchLowercase}
  \defaultfontfeatures[\rmfamily]{Ligatures=TeX,Scale=1}
\fi
\usepackage{lmodern}
\ifPDFTeX\else  
    % xetex/luatex font selection
\fi
% Use upquote if available, for straight quotes in verbatim environments
\IfFileExists{upquote.sty}{\usepackage{upquote}}{}
\IfFileExists{microtype.sty}{% use microtype if available
  \usepackage[]{microtype}
  \UseMicrotypeSet[protrusion]{basicmath} % disable protrusion for tt fonts
}{}
\makeatletter
\@ifundefined{KOMAClassName}{% if non-KOMA class
  \IfFileExists{parskip.sty}{%
    \usepackage{parskip}
  }{% else
    \setlength{\parindent}{0pt}
    \setlength{\parskip}{6pt plus 2pt minus 1pt}}
}{% if KOMA class
  \KOMAoptions{parskip=half}}
\makeatother
\usepackage{xcolor}
\setlength{\emergencystretch}{3em} % prevent overfull lines
\setcounter{secnumdepth}{-\maxdimen} % remove section numbering
% Make \paragraph and \subparagraph free-standing
\makeatletter
\ifx\paragraph\undefined\else
  \let\oldparagraph\paragraph
  \renewcommand{\paragraph}{
    \@ifstar
      \xxxParagraphStar
      \xxxParagraphNoStar
  }
  \newcommand{\xxxParagraphStar}[1]{\oldparagraph*{#1}\mbox{}}
  \newcommand{\xxxParagraphNoStar}[1]{\oldparagraph{#1}\mbox{}}
\fi
\ifx\subparagraph\undefined\else
  \let\oldsubparagraph\subparagraph
  \renewcommand{\subparagraph}{
    \@ifstar
      \xxxSubParagraphStar
      \xxxSubParagraphNoStar
  }
  \newcommand{\xxxSubParagraphStar}[1]{\oldsubparagraph*{#1}\mbox{}}
  \newcommand{\xxxSubParagraphNoStar}[1]{\oldsubparagraph{#1}\mbox{}}
\fi
\makeatother


\providecommand{\tightlist}{%
  \setlength{\itemsep}{0pt}\setlength{\parskip}{0pt}}\usepackage{longtable,booktabs,array}
\usepackage{calc} % for calculating minipage widths
% Correct order of tables after \paragraph or \subparagraph
\usepackage{etoolbox}
\makeatletter
\patchcmd\longtable{\par}{\if@noskipsec\mbox{}\fi\par}{}{}
\makeatother
% Allow footnotes in longtable head/foot
\IfFileExists{footnotehyper.sty}{\usepackage{footnotehyper}}{\usepackage{footnote}}
\makesavenoteenv{longtable}
\usepackage{graphicx}
\makeatletter
\newsavebox\pandoc@box
\newcommand*\pandocbounded[1]{% scales image to fit in text height/width
  \sbox\pandoc@box{#1}%
  \Gscale@div\@tempa{\textheight}{\dimexpr\ht\pandoc@box+\dp\pandoc@box\relax}%
  \Gscale@div\@tempb{\linewidth}{\wd\pandoc@box}%
  \ifdim\@tempb\p@<\@tempa\p@\let\@tempa\@tempb\fi% select the smaller of both
  \ifdim\@tempa\p@<\p@\scalebox{\@tempa}{\usebox\pandoc@box}%
  \else\usebox{\pandoc@box}%
  \fi%
}
% Set default figure placement to htbp
\def\fps@figure{htbp}
\makeatother

\usepackage{pdfpages}
\KOMAoption{captions}{tableheading}
\makeatletter
\@ifpackageloaded{caption}{}{\usepackage{caption}}
\AtBeginDocument{%
\ifdefined\contentsname
  \renewcommand*\contentsname{Table of contents}
\else
  \newcommand\contentsname{Table of contents}
\fi
\ifdefined\listfigurename
  \renewcommand*\listfigurename{List of Figures}
\else
  \newcommand\listfigurename{List of Figures}
\fi
\ifdefined\listtablename
  \renewcommand*\listtablename{List of Tables}
\else
  \newcommand\listtablename{List of Tables}
\fi
\ifdefined\figurename
  \renewcommand*\figurename{Figure}
\else
  \newcommand\figurename{Figure}
\fi
\ifdefined\tablename
  \renewcommand*\tablename{Table}
\else
  \newcommand\tablename{Table}
\fi
}
\@ifpackageloaded{float}{}{\usepackage{float}}
\floatstyle{ruled}
\@ifundefined{c@chapter}{\newfloat{codelisting}{h}{lop}}{\newfloat{codelisting}{h}{lop}[chapter]}
\floatname{codelisting}{Listing}
\newcommand*\listoflistings{\listof{codelisting}{List of Listings}}
\makeatother
\makeatletter
\makeatother
\makeatletter
\@ifpackageloaded{caption}{}{\usepackage{caption}}
\@ifpackageloaded{subcaption}{}{\usepackage{subcaption}}
\makeatother

\usepackage{bookmark}

\IfFileExists{xurl.sty}{\usepackage{xurl}}{} % add URL line breaks if available
\urlstyle{same} % disable monospaced font for URLs
\hypersetup{
  pdftitle={Projects},
  colorlinks=true,
  linkcolor={blue},
  filecolor={Maroon},
  citecolor={Blue},
  urlcolor={Blue},
  pdfcreator={LaTeX via pandoc}}


\title{Projects}
\author{}
\date{}

\begin{document}
\maketitle


\phantomsection\label{hero-heading}
\subsection{ESS 523C Labs}\label{ess-523c-labs}

\href{https://mrouin.github.io/csu-523c/lab-01-clean.html}{\textbf{Lab
1}} \textbf{- Exploring Streamflow Across the Colorado River Basin}

Retrieved 44 years of daily streamflow directly from the USGS NWIS API
using dataRetrieval. Joined gauge metadata using relational keys --- and
verified every join. Normalized by drainage area to enable fair
comparisons across the basin. Detected low-flow threshold exceedances
using rolling statistics for operational drought monitoring. Aggregated
by water year and season the natural temporal units of western
hydrology. Visualized basin-wide trends with a standardized anomaly
heatmap and a flow-weighted spatial center analysis. Modeled the
upstream-downstream relationship at the Compact's accounting point.
Evaluated model assumptions and identified the outlier year the model
couldn't predict

\href{https://mrouin.github.io/csu-523c/lab-02-clean.html}{\textbf{Lab
2}} \textbf{- The Geometry of Place: Spatial Analysis Skills for
Environmental Science}

Built spatial datasets from three different sources: a CRAN package
(tigris), a data package (rnaturalearth), and a downloaded CSV.
Transformed coordinate reference systems and selected a projection
appropriate for national-scale distance analysis. Manipulated geometry
types --- understanding when to use st\_union vs.~st\_combine, and why
casting to MULTILINESTRING is required before computing distances to
borders. Computed four sets of distances using st\_distance() and stored
them as labeled unit vectors. Visualized spatial patterns with
continuous color scales, gghighlight, and ggrepel. Built an interactive
web map using leaflet to present distance data to a general audience

\href{https://mrouin.github.io/csu-523c/lab-03-clean.html}{\textbf{Lab
3}} \textbf{- National Dam Inventory (NID): From Counts to Local
Clustering}

Tessellating surfaces to discretize space. Geometry simplification to
expedite expensive intersections. Writing functions to automate
repetitive reporting and mapping task. Point-in-polygon counts to
aggregate point data

\href{https://mrouin.github.io/csu-523c/lab-04-clean.html}{\textbf{Lab
4}} \textbf{- Palo Iowa Flooding}

Using remote sensing data we Identify an area of interest (AOI):
identified the temporal range of interest and the relevant images,
downloaded the images, analyzed the products.

\href{https://mrouin.github.io/csu-523c/lab-05-clean.html}{Lab 5}
\textbf{- Machine Learning in Hydrology}

In this lab we explored predictive modeling in hydrology using the
\texttt{tidymodels} framework and the CAMELS (Catchment Attributes and
Meteorology for Large-sample Studies) dataset. Specified a
\texttt{recipe}, dropped it into a \texttt{workflow}, and compared
multiple models via \texttt{workflow\_set} against cross-validated
resamples. Extracted \textbf{variable importance} and interpreted
predictor ranks through the \emph{geographic-proxy} lens. Evaluated with
\textbf{spatial cross-validation} and reported the honest gap between
random and spatial CV. \textbf{Tuned} hyperparameters with
\texttt{tune\_grid()} and a named grid, then finalized with
\texttt{last\_fit().}

\subsection{Cache la Poudre Stream
Metabolism}\label{cache-la-poudre-stream-metabolism}

\href{GitHub/mrouin/WR696\%20-\%20Project\%20Presentation.pdf}{Presentation}

Streams play an important role in communities, providing a source of
drinking water, infrastructure, recreation, and ecosystems for aquatic
life. Within Larimer County, Colorado, the Cache La Poudre River
primarily serves as a source for drinking water, agriculture, and
recreation for approximately 400,000 members within the watershed (CPRW
2026). Stream health within this watershed is affected by everyday
impacts, such as the agricultural and urban land uses, and by major
disturbance events such as the Cameron Peak Fire in 2020. These land use
and disturbance events drive nutrient and sediment loading into the
stream which can threaten the health of the stream and its ecosystem. To
understand these impacts, stream metabolism can be used to evaluate the
physical, biological, and chemical conditions of a stream and therefore
is a strong indicator of stream health (Reed et al.~2021).

The Cache la Poudre runs through varying land uses, from the
mountainous, rocky areas in the upper reaches to the urban and
agricultural uses in the middle and lower reaches. The variability of
this stream reach allows for a diverse system to analyze its metabolism.
The primary measurements that will be utilized when analyzing stream
metabolism in this paper are dissolved oxygen and temperature across
four sampling locations. This study aims to answer: How does the
metabolism of the Cache la Poudre River seasonally vary from the upper
to the lower reaches? To address this, I compare metabolism rates across
seasons at four sampling locations, starting below Chambers Reservoir
through to the CSU Environmental Learning Center. I also assess
longitudinal trends across sites and evaluate potential drivers of
observed patterns in stream metabolism.




\end{document}
